\cotefile{63}
\scanpage{57}
\begin{center}
CONGRUENCE HECKE OPERATORS

by

A.O.L.ATKIN
\end{center}

1. \handwritten{\uline}{Introduction}. \quad In \handwritten{\uwave}{(2)} the au\deleted{l}thor has given an account of some results on congruences for modular forms based on real time order of discovery,ind\deleted{\unsure{i}}\unsure{e|i}spensably aided by a computer. Here we attempt(with a certain amount of hindsight)to start at the other end,and consider what kind of relevant congruence properties might be expected \handwritten{\uline}{a} \handwritten{\uline}{priori} ; our conclusion is,that something like the results of \handwritten{\uwave}{(2)} are the most that can be hoped for.

We give no proofs in this paper. The historical aspect,and a precise account of what we can prove and what is conjecture(with overwhelming evidence)is in \handwritten{\uwave}{(2)}; detailed proofs will appear elsewhere. Our purpose here is to indicate that some \reversemarginpar\margin{\handwritten{German\\ `pee'}}\normalmarginpar``$\handwritten{\underline}{\handwritten{\mathfrak{p}}}$-adic'' theory of Hecke operators exists analogous to the exact classical theory. We sketch this classical theory of Hecke operators with a view to our analogy,so that the approach is mildly eccentric. The reader can find more satisfactory accounts in Gunning\handwritten{\uwave}{(4)},Lehner\handwritten{\uwave}{(7)}, or Rankin\handwritten{\uwave}{(8)}.

We observe finally that,while we only present our conjectures and theorems for the Fourier coefficients of modular functions(i.e.,forms of dimension $0$)and the subgroup $\handwritten{\underline}{\handwritten{\Gamma}}_{\handwritten{0}}(\handwritten{\underline}{q})$\margin{\handwritten{Cap.\\ gamma}} for $\handwritten{\underline}{q}$ prime,we have in fact a good deal of further evidence for forms of various dimensions(including half-integral)and other subgroups. In particular,for the reader whose interests lie in the ``romantic'' tradition of Ramanujan's famous formulae for the partition function $\handwritten{\underline}{p(n)}$:
\[
\begin{aligned}
\handwritten{\underline}{p(5n+4)}  &\handwritten{\equiv} 0 \qquad (\mathrm{mod}\ 5)  &,\\
\handwritten{\underline}{p(7n+5)}  &\handwritten{\equiv} 0 \qquad (\mathrm{mod}\ 7)  &,\\
\handwritten{\underline}{p(11n+6)} &\handwritten{\equiv} 0 \qquad (\mathrm{mod}\ 11) &,
\end{aligned}
\]
we can offer such results as
\[
\handwritten{\underline}{p(206839n+2623)}\ \handwritten{\equiv}\ 0 \qquad (\mathrm{mod}\ 17).
\]
However we can only prove this by means of the multiplicative congruence properties which we consider below.

% archive sheet number: 56

\scanpage{58}
\noindent\hfill atkin page 2\mspage{2}\par
\begingroup
\csname tagsleft@true\endcsname

2. \handwritten{\uline}{The} \handwritten{\uline}{full} \handwritten{\uline}{modular} \handwritten{\uline}{group}. We consider the modular group $\handwritten{\underline}{\handwritten{\Gamma}}$ consisting of all linear fractional transformations
\[
z\ \unsure{\rightarrow|-}\ Vz=(az+b)/(cz+d)
\quad \text{(with } \handwritten{\underline}{a,b,c,d},\text{rational integers and }\handwritten{\underline}{ad-bc\mathrel{\handwritten{=}}}1)
\]
acting on the upper half $\handwritten{\underline}{z}$-plane $\handwritten{\underline}{H}$,and write
\[
\handwritten{\underline}{x}=\handwritten{\underline}{e}^{2\handwritten{\underline}{\handwritten{\pi}}\handwritten{\underline}{iz}} \qquad \unsure{.}
(\,|x|<1\tikz[remember picture,baseline]{\node[inner sep=0pt](punctfiftyeight){$\deleted{\cdot}\mathpunct{,}$};}\quad \mathrm{Im}\ \handwritten{\underline}{z}>0\,)\ .
\]
\reversemarginpar\margin{\handwritten{\unsure{.} pi}}\normalmarginpar
\margin{\handwritten{\tikz[remember picture,baseline]{\node[inner sep=0pt](delefiftyeight){\tikz[baseline=-.4ex,x=.85em,y=.85em]{\draw (.03,-.7) .. controls (.36,-.06) and (.83,.64) .. (.59,.63) .. controls (.18,.63) and (.16,.06) .. (.45,.22) .. controls (.61,.35) and (.32,.56) .. (.46,.68) -- (.93,1.12);}};}\tikz[remember picture,overlay]{\draw[->] (punctfiftyeight.north east) to[bend left=6] (delefiftyeight.west);}}}

We define a \handwritten{\uline}{modular} \handwritten{\uline}{form} $\handwritten{\underline}{F(z)}$ of dimension $\handwritten{\underline}{k}$(which we take to be an even integer throughout this paper) by
\[
\handwritten{\underline}{F(Vz)=(cz+d)^{-k}F(z)\quad \text{for all } V \handwritten{\in}\handwritten{\underline}{\handwritten{\Gamma}}}\quad ,
\tag*{(M1)}
\]
\[
\parbox[t]{.82\linewidth}{\raggedright \handwritten{\uline}{$F(z)$ is meromorphic in the fundamental domain $D=I+B$ of $\handwritten{\underline}{\handwritten{\Gamma}}$} ,}
\tag*{(M2)}
\]
\handwritten{\uline}{where $I$ is that part of $D$ inside $H$,and $B$(the cusps of $D$) is that part of $D$ on the boundary of $H$; here we may take $B$ as the single point $z=i\handwritten{\infty}$.}

The conditions (M1) and (M2) seem insufficient to produce much number theory,and we therefore define further restrictions as follows:
\[
\text{an \handwritten{\uline}{entire} form is regular in } I\quad ,
\tag*{(M3a)}
\]
\[
\text{a \handwritten{\uline}{regular} form is regular in } D=I+B\quad ,
\tag*{(M3b)}
\]
\[
\text{a \handwritten{\uline}{cuspform} is regular in } I \text{ and zero in } B.
\tag*{(M3c)}
\]
Clearly a cuspform is a regular form,and a regular form is an entire form. We denote the vector spaces of entire,regular,and cusp-,forms(for the given dimension $\handwritten{\underline}{k}$) by $\handwritten{\underline}{C},\handwritten{\underline}{C}^{\handwritten{+}}$,and $\handwritten{\underline}{C}^{\handwritten{0}}\handwritten{,}$respectively. The last two \tikz[baseline=(old.base)]{\node[inner sep=0pt](old){\deleted{\unsure{aure|aare}}};\draw (old.center) ellipse [x radius=1em,y radius=.65em];}\inserted{\handwritten{\tikz[baseline=(new.base)]{\node[inner sep=0pt](new){are};\draw (new.center) ellipse [x radius=.9em,y radius=.65em];}}} finite-dimensional($\handwritten{\underline}{C}$ is not). Since $\handwritten{\underline}{C}^{\handwritten{+}}$ is the direct sum of $\handwritten{\underline}{C}^{\handwritten{0}}$ and a well-defined Eisenstein series,we ignore $\handwritten{\underline}{C}^{\handwritten{+}}$ from now on.

By the \handwritten{\uline}{Fourier} \handwritten{\uline}{series} (FS)\unsure{-} of a modular form we understand its expansion in powers of $\handwritten{\underline}{x}$,the local variable at the cu\deleted{p}sp $\handwritten{\underline}{z}=\handwritten{\underline}{i}\handwritten{\infty}$ $\handwritten{\unsure{\vdots|\mathord{:}}}$
\[
F(z)=F(x)=\handwritten{\sum_{n=N}^{\infty}}a(n)\,x^{n}\qquad (N>-\handwritten{\infty})\qquad ,
\]
and we call \tikz[remember picture,baseline=(oldcoeff.base)]{\node[inner sep=0pt](oldcoeff){\deleted{\illegible{1}}};}\reversemarginpar\margin{\handwritten{\tikz[remember picture,baseline]{\node[inner sep=0pt](deletwo){\tikz[baseline=-.4ex,x=.85em,y=.85em]{\draw (.03,-.7) .. controls (.36,-.06) and (.83,.64) .. (.59,.63) .. controls (.18,.63) and (.16,.06) .. (.45,.22) .. controls (.61,.35) and (.32,.56) .. (.46,.68) -- (.93,1.12);}};}\tikz[remember picture,overlay]{\draw (oldcoeff.north) to[bend right=7] (deletwo.east);}}}\normalmarginpar $\handwritten{\underline}{a(n)}$ the Fourier coefficients. Roughly speaking, the number theory associated with modular theory arises from the Fourier coefficients of entire modular forms. Indeed,prior to Ramanujan and his work on congruences and(with Hardy)asymptotic formulae for $\handwritten{\underline}{p(n)}$,one could read ``regular''for ``entire'' in the last sentence. The author does not know of any number theory associated with modular forms which have poles inside $\handwritten{\underline}{H}$.
\endgroup

% archive sheet number: 57

\scanpage{59}
\noindent\hfill\unsure{atkin page 3}\mspage{3}\par
\begingroup
\csname tagsleft@true\endcsname

We now define the Hecke operator $\handwritten{\underline}{T_p}$ for prime $\handwritten{\underline}{p}$ acting on a modular form $\handwritten{\underline}{\deleted{\unsure{E|F}}\handwritten{F}(z)}$ of dimension $\handwritten{\underline}{k}$ by
\[
T_p\,F(z)=p^{-1}\handwritten{\sum_{\underline{\lambda}=0}^{p-1}F\left(\frac{z+\underline{\lambda}}{p}\right)}
+p^{-k-1}F(pz),
\]
\margin{\handwritten{lambda}}
or equivalently in terms of the FS:
\[
T_p\,F(x)=\handwritten{\sum_n}\handwritten{\Bigl\{}a(np)+\unsure{\cdot}p^{-k-1}\,a(n/p)\handwritten{\Bigr\}}\,x^n\qquad ,
\]
where $\handwritten{\underline}{a}(\handwritten{\underline}{\alpha})$ is defined as zero if $\handwritten{\underline}{\alpha}$ is non-integral or $\handwritten{\underline}{\alpha}<\handwritten{\underline}{N}$.\margin{\handwritten{alpha}} It is then a matter of formal algebra to establish the two following properties:
\[
\text{\handwritten{\uline}{$T_p$ maps $C$ into $C$,and $C^{\handwritten{0}}$ into $C^{\handwritten{0}}$}}\ ,
\tag*{(H1)}
\]
\[
\text{\handwritten{\uline}{All the $T_p$ commute.}}
\tag*{(H2)}
\]
The further development of the classical theory now requires Petersson's scalar(inner)product
\[
(F,G)=\handwritten{\iint_{D}}\,F(z)\,\handwritten{\overline}{G(z)}\,\handwritten{\underline{\eta}^{-k-2}\,d\underline{\xi}\,d\underline{\eta}}
\qquad (z=\handwritten{\underline{\xi}+i\underline{\eta}})\ ,
\]
\margin{\handwritten{\unsure{gsi|ksi}\\ eta}}
which co\handwritten{n}verges provided $\handwritten{\underline}{F,G}\handwritten{\in}\handwritten{\underline}{C}^{\handwritten{+}}$ and at least one of $\handwritten{\underline}{F,G}\handwritten{\in}\handwritten{\underline}{C}^{\handwritten{0}}$, and makes $\handwritten{\underline}{C}^{\handwritten{0}}$ into a Hilbert space. We then have
\[
\text{\handwritten{\uline}{$T_p$ is an hermitian operator on $C^{\handwritten{0}}$,i.e.}}, 
\tag*{(H3)}
\]
\[
(T_pF,G)=(F,T_pG)\ ,
\]
and we have already stated
\margin{\handwritten{\tikz[x=.6em,y=.6em]{\draw[very thin] (.3,2.1) .. controls (.95,1.7) and (.9,1.4) .. (.17,1.51) .. controls (-.5,1.72) and (-.27,1.12) .. (.25,1.33) .. controls (.75,1.61) and (.5,.52) .. (1,-.6);}}}
\[
\text{\handwritten{\uline}{The Hilbert space $C^{\handwritten{0}}$ is finite-dimensional}.}
\tag*{(H4)}
\]
The properties (H1) through (H4) suffice to establish the following

\handwritten{\uline}{Theorem} 1. (\handwritten{\uline}{Hecke-Petersson}).

\handwritten{\uline}{The space of cuspforms of dimension $k$ on the full modular group has a basis $F_i(z)$ of forms which are simultaneous eigenforms}*\margin{\handwritten{*footnote\\ on \unsure{cover}\\ page}} \handwritten{\uline}{of all the operators $T_p$} . \handwritten{\uline}{In terms of the Fourier series} $\handwritten{\sum\limits_{n=1}^{\infty}}\,a(n)\,x^{n}$ \handwritten{\uline}{of a typical} $F_i(z)$ \handwritten{\uline}{normalized by} $a(1)=1$ \handwritten{\uline}{this is equivalent to the ``multiplicative'' relation}
\csname tagsleft@false\endcsname
\[
\handwritten{\underline}{a(np)-a(n)a(p)+p^{-k-1}a(n/p)=0}\quad
(\text{\handwritten{\uline}{all prime} }p,\text{\handwritten{\uline}{all integral} }n)\ .
\tag*{( 1 )\,\handwritten{\#}}
\]
From (1) and the values of $\handwritten{\underline}{a(p)}$ for prime $\handwritten{\underline}{p}$ the values of $\handwritten{\underline}{a(n)}$ are completely determined; in particular $\handwritten{\underline}{a(mn)=a(m)a(n)}$ if $\handwritten{\underline}{(m,n)=1}$. For this reason we have not defined the operators $\handwritten{\underline}{T_m}$ for composite $\handwritten{\underline}{m}$ which form a ring,but merely the \deleted{group}\inserted{\handwritten{\unsure{multiplicative|multiplicatif} \uline{semi} \unsure{group|grup}!}} of $\handwritten{\underline}{T_p}$ ,since the number theory(if not the algebra) is completely described by the equation (1). One can infer Theorem 1 in any specific case with a much weaker form of(H3),and indeed without the scalar product; it suffices to have (H1),(H2),(H4),together with
\csname tagsleft@true\endcsname
\[
\begin{aligned}
&\text{\handwritten{\uline}{At least one} } T_p\ \text{(\handwritten{\uline}{considered as a matrix acting on some basis}}\\
&\text{\unsure{for|of} } C^{\handwritten{0}}\text{ ) \unsure{has distinct} eigenvalues.}
\end{aligned}
\tag*{(H3\unsure{*|$\cdot$})}
\]
\endgroup

% archive sheet number: 58

\scanpage{60}
\begin{center}atkin page 4\mspage{4}\end{center}

If we now consider the possibility of extending some form of Theorem 1 to $\handwritten{\underline}{C}$(rather than $\handwritten{\underline}{C}^{0}$ )then there is no problem with (H1) and (H2). However\deleted{\unsure{,}} (H4)\handwritten{,}which is essential\handwritten{,}breaks down since the only finite-dimensional subspaces of $\handwritten{\underline}{C}$ which are mapped into themselves by all the $\handwritten{\underline}{T_p}$ are subspaces of $\handwritten{\underline}{C}^{\handwritten{+}}$. Indeed any modular form with a pole\handwritten{\tikz[baseline,overlay,x=1em,y=1em]{\draw[very thin,gray] (-.9,-.35) .. controls (-.7,-.63) and (.85,-.8) .. (1.3,-1) .. controls (1.55,-1.1) and (1.39,-.59) .. (1.58,-.66) .. controls (1.76,-.79) and (1.63,-.92) .. (1.48,-1.06);}} gives rise to an infinite-dimensional subspace by the repeated action of only one $\handwritten{\underline}{T_p}$ .

The essential distinction,then,lies between $\handwritten{\underline}{C}$ and $\handwritten{\underline}{C}^{\handwritten{0}}$. It so happens that $\handwritten{\underline}{C}^{\handwritten{0}}$ is null for $\handwritten{\underline}{k}\handwritten{\geqslant}0$ and there is thus an association between Theorem 1 and forms of negative dimension. In fact such theory as we are able to develop for $\handwritten{\underline}{C}$ is viable for forms of any dimension,but the author has mainly pursued it for positive and zero dimension where nothing else is available.

We now choose the simplest of all modular forms,in fact the modular function($\handwritten{\underline}{k}=0$) $\handwritten{\underline}{j(z)}$ which is the Hauptmodul of $\handwritten{\underline}{\handwritten{\Gamma}}$ defined by
\[
\begin{aligned}
j(z)&=\handwritten{\left\{1+240\sum_{n=1}^{\infty}\underline{\sigma}_{3}(n)\,x^n\right\}^{3}\Bigm/ x\prod_{r=1}^{\infty}(1-x^{r})^{24}}\\
&=x^{-1}+744+\handwritten{\sum_{n=1}^{\infty}}c(n)\,x^n,
\end{aligned}
\]
\reversemarginpar\margin{\handwritten{sigma}}\normalmarginpar
and consider what kind of property
\[
a(1)=1,\ a(np)-a(n)a(p)+p^{-1}a(n/p)=0\quad (\text{all prime }p,\text{all integral }n),
\tag*{(2)}
\]
can exist for $\handwritten{\underline}{a(n)}$ associated in some way with $\handwritten{\underline}{c(n)}$. Nothing identical is possible,for the reasons already given. However (2) makes sense if the $\handwritten{\underline}{a(n)}$ are considered not\inserted{necessarily} as complex numbers but as elements of any commutative ring with identity; since in fact the $\handwritten{\underline}{a(n)}$ of Theorem 1 are algebraic integers(as shown by Rankin and Rushforth\handwritten{\uwave}{(9)}) a natural choice of ring is $\handwritten{\underline}{Z/q^{\handwritten{\underline}{\alpha}}Z}$\margin{\handwritten{alpha}} ,that is, the integers modulo $\handwritten{\underline}{q}^{\handwritten{\underline}{\alpha}}$,for some fixed prime $\handwritten{\underline}{q}$. We cannot expect some \handwritten{\uline}{fixed} $\handwritten{\underline}{a(n)}$ to be involved for all $\handwritten{\underline}{q}$,or even for fixed $\handwritten{\underline}{q}$ and all $\handwritten{\underline}{\alpha}$,for then (2) would be identical again. Thus we may expect $\handwritten{\underline}{a(n)}$ to depend on $\handwritten{\underline}{q}$ and $\handwritten{\underline}{\alpha}$ . At this point we \deleted{\illegible{1}} state some actual results,which indicate that something of this kind actually happens; for $\handwritten{\underline}{q}=2,3,5,7$,and $13$ we can prove for all $\handwritten{\underline}{\alpha}$,and otherwise\unsure{.}for a few fixed $\unsure{a|\alpha}$ the following

% archive sheet number: 59

\scanpage{61}

\handwritten{\uline}{Conjecture} 1.

\handwritten{\tikz[baseline,overlay,x=1em,y=1em]{\draw[very thin,gray] (-1.55,-1.12) -- (-.72,-.51) -- (-.48,-.77) (-.76,-.44) -- (-.3,-.49);}}\handwritten{\uline}{Let} $q\handwritten{\leqslant}23$ be \handwritten{\uline}{prime},\handwritten{\uline}{and} $p$ \inserted{\handwritten{be}} a \handwritten{\uline}{prime other than} $q$,\handwritten{\uline}{and let $\handwritten{\underline}{\alpha}$\inserted{and $n$} be any positive integers} . \handwritten{\uline}{Let also}
\[
a(n)=c(q^{\handwritten{\underline}{\alpha}}n)/c(q^{\handwritten{\underline}{\alpha}})\ .
\]
\handwritten{\uline}{Then}
\[
\begin{aligned}
a(np)-a(n)a(p)+p^{-1}a(n/p)&\handwritten{\equiv}0 \qquad (\mathrm{mod}\ q^{\handwritten{\underline}{\alpha}})\qquad ,\\
a(nq)-a(n)a(q)&\handwritten{\equiv}0 \qquad (\mathrm{mod}\ q^{\handwritten{\underline}{\alpha}}).
\end{aligned}
\]
\margin{\handwritten{\#}}

In fact for $\handwritten{\underline}{q}\handwritten{\leqslant}5$ the modulus of the congruences can be improved to $q^{\handwritten{\underline}{\theta}}$\reversemarginpar\margin{\handwritten{theta}}\normalmarginpar with $\handwritten{\underline}{\theta}>\handwritten{\underline}{\alpha}$ . The reader familiar with the Ramanujan-type congruences(Lehner\handwritten{\uwave}{(5)},\handwritten{\uwave}{(6)},Atkin\handwritten{\uwave}{(1)})
\[
c(q^{\handwritten{\underline}{\alpha}}n)\handwritten{\equiv}0\quad(\mathrm{mod}\ q^{\handwritten{\underline}{\beta}})
\tag*{(3)}
\]
for $\handwritten{\underline}{q}=2,3,5,7,11$ and $\handwritten{\underline}{\beta}=3\handwritten{\underline}{\alpha}+8,2\handwritten{\underline}{\alpha}+3,\handwritten{\underline}{\alpha}+1,\handwritten{\underline}{\alpha},\handwritten{\underline}{\alpha}$\reversemarginpar\margin{\handwritten{beta}}\normalmarginpar respectively may query the definition of $\handwritten{\underline}{a(n)}$,but in fact for $\handwritten{\underline}{n}=1$ the result of (3) is best possible(i.\handwritten{e}. $\handwritten{\underline}{q^{\handwritten{\underline}{\beta}}\handwritten{\parallel}c(q^{\handwritten{\underline}{\alpha}})}$),and the division is therefore valid in any ring $\handwritten{\underline}{Z/q^{\handwritten{\underline}{\theta}}Z}$ . Next we state

\handwritten{\uline}{Conjecture} 2. \quad \handwritten{\uline}{Let} $\handwritten{\underline}{q=29\ \text{or}}$ $31$,\inserted{\handwritten{and}} $p$ \handwritten{\uline}{be a prime other than} $q$,\handwritten{\uline}{and $\handwritten{\underline}{\alpha}$ and $n$ be any positive integers}. \handwritten{\uline}{Then}
\[
c(q^{\handwritten{\underline}{\alpha}}n)/c(q^{\handwritten{\underline}{\alpha}})=\handwritten{\underline}{\handwritten{\lambda}}a(n)+\handwritten{\underline}{\handwritten{\bar{\lambda}}}\,\handwritten{\overline}{a(n)}\ ,
\]
\handwritten{\uline}{where}$\handwritten{\underline}{\handwritten{\lambda}},a(n)$ \handwritten{\uline}{and their conjugates}$\handwritten{\underline{\bar{\lambda}}},\handwritten{\overline}{a(n)}$ \handwritten{\uline}{lie in a quadratic extension} $E$ of $Z/q^{\handwritten{\underline}{\alpha}}Z$ \handwritten{\uline}{and}
\[
\begin{aligned}
a(np)-a(n)a(p)+p^{-1}a(n/p)&=0\ ,\\
a(nq)-a(n)a(q)&=0.
\end{aligned}
\]
\margin{\handwritten{\#}}
(We have written ``$=$'' rather than ``$\handwritten{\equiv}$'' in $E$ \ ). The situation for $\handwritten{\underline}{q}\handwritten{\geqslant}37$ appears to be that something like \deleted{c}\handwritten{C}onjecture 2 is valid; one finds that $\handwritten{\underline}{c(q^{\handwritten{\underline}{\alpha}}n)/c(q^{\handwritten{\underline}{\alpha}})}$ is a linear combination of ``congruence multiplicative'' coefficients lying in extensions of $\handwritten{\underline}{Z/q^{\handwritten{\underline}{\alpha}}Z}$ and the ``dimension'' of the``congruence space'' appears to be the integral part of $\handwritten{\underline}{q}/12$ (this is somewhat speculative) in the sense that the ``dimension'' is 1 in conjecture 1 and 2 in Conjecture 2.

We have so far considered the coefficient $\handwritten{\underline}{c(n)}$ on the full modular group. H\inserted{\handwritten{o}}wever the results quoted above give the prime $q$ a sp ecial role,\handwritten{\tikz[baseline,overlay,x=1em,y=1em]{\draw[very thin] (1.16,1.03) .. controls (1.29,1.31) and (1.46,1.16) .. (1.62,.69);}}both as to the definition of $\handwritten{\underline}{a(n)}$ and the multiplicative property; also $\handwritten{\sum\limits_{n=1}^{\infty}}c(q^{\handwritten{\underline}{\alpha}}n)\,x^{n}$ is (for all $\handwritten{\underline}{\alpha}$ ) a function on the subgroup $\handwritten{\underline}{\handwritten{\Gamma}}_{\handwritten{0}}(\handwritten{\underline}{q})$ defined by $\handwritten{\underline}{c\handwritten{\equiv}0\ (\mathrm{mod}\ q)}$ in

$z\ \unsure{\rightarrow|-}\ Vz=(az+b)/(cz+d)$. It therefore seems that these results are naturally associated with $\handwritten{\underline}{\handwritten{\Gamma}}_{\handwritten{0}}(q)$ and we now consider this

% archive sheet number: 60

\scanpage{62}
in more detail.

3. \handwritten{\uline}{The} \handwritten{\uline}{subgroup} $\handwritten{\underline}{\handwritten{\Gamma}}_{\handwritten{0}}(q)$ \handwritten{\uline}{for} $q$ \handwritten{\uline}{prime}.

The theory of \handwritten{\S}2 carries through with some minor changes,and the replacement of the operator $\handwritten{\underline}{T_q}$ by the operator $\handwritten{\underline}{U_q}$ defined by
\[
 U_q F(z)=q^{-1}\mathop{\handwritten{\sum}}\limits_{\handwritten{\underline}{\handwritten{\lambda}}=\handwritten{0}}^{\handwritten{q-1}} \handwritten{F\left(\frac{z+\underline{\lambda}}{q}\right)}=\mathop{\handwritten{\sum}}\limits_{\handwritten{n}}a(nq)\,x^n\,.
\]
(M1) need only hold for $V \handwritten{\in}\handwritten{\underline}{\handwritten{\Gamma}}_{\handwritten{0}}(q)$,and in (M2) the \handwritten{\uline}{two} cusps $\handwritten{\underline}{B}$ can be taken as $z=0$ and $z=i\handwritten{\infty}$. The existence of another cusp enables us to have\inserted{entire} functions which are as it were ``cusp'' functions with a zero at $\handwritten{\underline}{z}=i\handwritten{\infty}$ by virtue of having a pole at $\handwritten{\underline}{z}=0$ . With regard to the classical theory it is convenient to discard the cuspforms $\handwritten{\underline}{G(z)}$ and $\handwritten{\underline}{G(qz)}$ on $\handwritten{\underline}{\handwritten{\Gamma}}_{\handwritten{0}}(\handwritten{\underline}{q})$ where $\handwritten{\underline}{G(z)}$ is a cuspform of the same dimension on the full modular group(and so a fortiori on $\handwritten{\underline}{\handwritten{\Gamma}}_{\handwritten{0}}(q)$ ),and to consider $C^0$ as the space of ``newforms'' which are the remaining cuspforms on $\handwritten{\underline}{\handwritten{\Gamma}}_{\handwritten{0}}(\handwritten{\underline}{q})$ of dimension $\handwritten{\underline}{k}$. We then have

\deleted{\unsure{Theorem}}
\begingroup\csname tagsleft@true\endcsname
\[
 \handwritten{\underline}{T_p\text{ and }U_q\text{ map }C\text{ into }C\text{ and }C^{\handwritten{0}}\text{ into }C^{\handwritten{0}}}\,, \tag*{(H1q)}
\]
\[
 \handwritten{\underline}{\text{All the } \deleted{T\unsure{s|p}}\handwritten{T_p}\text{ commute,and commute with }U}_{q}\,, \tag*{(H2q)}
\]
\[
 \handwritten{\underline}{T_p\text{ is an hermitian operator on }C^{\handwritten{0}}\text{ (note }U_q\text{ is not)}}\,, \tag*{(H3q)}
\]
\[
 \handwritten{\underline}{C^{\handwritten{0}}\text{ is finite dimensional}}. \tag*{(H4q)}
\]
\endgroup

This leads (with a little extra work to deal with $U_q$) to

\handwritten{\uline}{Theorem} 2. \handwritten{\uline}{The space of new cuspforms on} $\handwritten{\underline}{\handwritten{\Gamma}}_{\handwritten{0}}(q)$ \handwritten{\uline}{has a basis} $F_i(z)$ \handwritten{\uline}{of forms} \handwritten{\uline}{which are simultaneous eigenforms} of \handwritten{\uline}{all the operators} $T_p$ \handwritten{\uline}{and of} $U_q$ . In \handwritten{\uline}{terms of the Fourier series} $\mathop{\handwritten{\sum}}\limits_{\handwritten{n=1}}^{\handwritten{\infty}}a(n)\,x^n$ of \handwritten{\uline}{a typica}\deleted{\unsure{l|o}}l $F_i(z)$ \handwritten{\uline}{normalized by} $a(1)=1$ this \handwritten{\uline}{is equivalent to the multiplicative} \handwritten{\uline}{relations}
\[
 \begin{aligned}
 a(np)-a(n)a(p)+p^{-k-1}a(n/p)&=0,\qquad (\text{\handwritten{\uline}{all prime} }p\handwritten{\ne}q,\text{\handwritten{\uline}{all integral} } n),\\
 a(nq)-a(n)a(q)&=0. 
 \end{aligned}
\]
\margin{\handwritten{\#}}
(In fact we have,in the case of negative $k$ where Theorem 2 is non-trivial,that $a(q)=\mathbin{\handwritten{\pm}}q^{-k/2\,-1}$. The theory on $\handwritten{\underline}{\handwritten{\Gamma}}_{\handwritten{0}}(m)$ for a general $\handwritten{\underline}{m}$ will appear in a paper by Joseph Lehner and the author). We now state\deleted{\unsure{a form|form}}\inserted{an extension} of Conjecture 1 which is more closely analogous to Theorem 2.
% archive sheet number: 61

\scanpage{63}
\begin{center}atkin page \handwritten{7}\end{center}\mspage{7}

\handwritten{\uline}{Conjecture} 3. \handwritten{\uline}{Let} $q\handwritten{\leqslant}23$ \handwritten{\uline}{be prime},\handwritten{\uline}{and $p$ be a prime other than} $q$, \handwritten{\uline}{and $n$ be a positive integer}. \handwritten{\uline}{Then for integral} $\handwritten{\underline}{\alpha}\handwritten{\geqslant}1$ \handwritten{\uline}{there exists a sequence of entire functions} $F_{\alpha}(z)=F_{\alpha}(x)$ \handwritten{\uline}{on} $\handwritten{\underline}{\handwritten{\Gamma}}_{\handwritten{0}}(q)$ \handwritten{\uline}{which are zero} at $z=i\handwritten{\infty}$ \handwritten{\uline}{(and hence have poles at $z=0$),such that if the FS of} $F_{\handwritten{\underline}{\alpha}}(x)$ \handwritten{\uline}{is given by}
\[
 F_{\alpha}(x)=x+\mathop{\handwritten{\sum}}\limits_{\handwritten{n=1}}^{\handwritten{\infty}}a_{\handwritten{\underline}{\alpha}}(n)\,x^n
 \qquad(\text{the }a_{\alpha}(n)\text{ are integral}),
\]
\handwritten{\uline}{then}
\[
 \begin{aligned}
 \text{(i\handwritten{\tikz[baseline,overlay,x=1em,y=1em]{\draw[line width=.08em] (0,.23) -- (.30,.23);}})}\quad &a_{\handwritten{\underline}{\alpha}}(np)-a_{\handwritten{\underline}{\alpha}}(n)a_{\handwritten{\underline}{\alpha}}(p)+p^{-1}a_{\handwritten{\underline}{\alpha}}(n/p)\handwritten{\equiv}0 &(\mathrm{mod}\ q^{\alpha})\,,\\
 \text{(ii)}\quad &a_{\handwritten{\underline}{\alpha}}(nq)-a_{\handwritten{\underline}{\alpha}}(n)a_{\handwritten{\underline}{\alpha}}(q)\handwritten{\equiv}0 &(\mathrm{mod}\ q^{\alpha})\,,\\
 \text{(iii)}\quad &a_{\handwritten{\underline}{\alpha}+1}(n)\handwritten{\equiv}a_{\handwritten{\underline}{\alpha}}(n) &(\mathrm{mod}\ q^{\alpha})\,,
 \end{aligned}
\]
(iv) \handwritten{\uline}{any entire function} $F(x)$ \handwritten{\uline}{on} $\handwritten{\underline}{\handwritten{\Gamma}}_{\handwritten{0}}(q)$ \handwritten{\uline}{with integral} \handwritten{\uline}{FS is such} \handwritten{\uline}{that for} \handwritten{\uline}{some} $\beta$ \handwritten{\uline}{we have} \inserted{\handwritten{for all $\alpha\geqslant1$}}
\[
 U_q^{\handwritten{\underline}{\alpha+\beta}} F(x)\handwritten{\equiv}\handwritten{k+}k_{\handwritten{\underline}{\alpha\beta},} F_{\handwritten{\underline}{\alpha}}(z)
 \qquad(\mathrm{mod}\ q^{\alpha})\,,
\]
\handwritten{\uline}{where}\inserted{\handwritten{$k$ and}} $k_{\handwritten{\underline}{\alpha\beta}}$ \deleted{\unsure{is}}\handwritten{are} \handwritten{\uline}{constant}\handwritten{s.}\margin{\handwritten{\#}}

It seems possible that such an $F_{\handwritten{\underline}{\alpha}}(x)$ is unique but the author can see no approach to prove this. It is clearly attractive,in view of (iii),to replace the sequence of functions $F_{\handwritten{\underline}{\alpha}}(x)$ by a single function $\handwritten{\underline}{F(x)}$ with $\handwritten{\underline}{q}$-adic Fourier coefficients. However the author's only methods of proof are based on classical modular theory,and for each fixed $\handwritten{\underline}{\alpha}$ we do have a genuine function on $\handwritten{\underline}{\handwritten{\Gamma}}_{\handwritten{0}}(\handwritten{\underline}{q})$ with a pole of finite order at $\handwritten{\underline}{z}=0$;one can prove the congruence properties via the exact properties. One may hope for some deeper method of proof,which in view of the exciting modern applications of algebraic geometry to modular theory is perhaps not wildly optimistic.

We may sum up the conclusions of Conjecture 3 as follows. Given ``any'' entire function on $\handwritten{\underline}{\handwritten{\Gamma}}_{\handwritten{0}}(\handwritten{\underline}{q})$, then repeated application of $\handwritten{\underline}{U_q}$ leads to a series of Fourier coefficients which become multi\deleted{\unsure{p|i}}\handwritten{p}licative modulo increasing powers of $\handwritten{\underline}{q}$. It is true that we have made the condition that the original FS be integral,but this is not an essential restriction,since it seems likely that any\inserted{entire} form on any \handwritten{\uline}{congruence} subgroup of the modular group is a finite linear combination of forms with integral FS(and this is certainly \deleted{\unsure{t|i}}true for functions on $\handwritten{\underline}{\handwritten{\Gamma}}_{\handwritten{0}}(\handwritten{\underline}{q})$ with $\handwritten{\underline}{q}\handwritten{\leqslant}23$ ).
\par\noindent\hspace*{.25\linewidth}\handwritten{\tikz[baseline,x=1em,y=1em]{\fill (-.58,-.12) circle (.025); \draw[line width=.045em] (.15,.51) .. controls (.08,.38) and (.13,.36) .. (.17,.25) .. controls (.28,.13) and (.01,.12) .. (.02,.05) .. controls (.05,-.02) and (.18,.01) .. (.2,0);}}
% archive sheet number: 62

\scanpage{64}
\begin{flushright}atkin page 8\end{flushright}\mspage{8}

Finally we may extend Conjecture 2 in the same way as Conjecture \handwritten{3} extends Conjecture 1. Our proofs(where available) have two difficulties; the first,and most crucial,is to obtain the analogy of (H4q) and show that congruentially we can\inserted{for all $\handwritten{\underline}{\alpha}$} obtain a finite-dimensional subspace of $\handwritten{\underline}{C}$,which we can only do when $\handwritten{\underline}{\handwritten{\Gamma}}_{\handwritten{0}}(\handwritten{\underline}{q})$ has genus zero by the methods of Atkin and O'Brien\handwritten{\uwave}{(3)},where the case $\handwritten{\underline}{q}=13$ of Conjecture 1 is stated. For the remaining values of $\handwritten{\underline}{q}$ we are limited to a few fixed values of $\handwritten{\underline}{\alpha}$. The second difficulty,which only arises for $\handwritten{\underline}{q}\handwritten{\geqslant}29$ when the ``congruence space'' has dimension greater than one,is to obtain a criterion for simultaneous diagonalization such as (H3q). In fact we can use an (H3q*) analogous to (H3*) in \handwritten{\S}2 to the effect that at least one matrix has distinct eigenvalues modulo $\handwritten{\underline}{q}$(it will a fortiori then have distinct eigenvalues modulo $\handwritten{\underline}{q}^{\handwritten{\underline}{\alpha}}$). As indicated in \handwritten{\S}2,there is no problem with (H1q) and (H2q) which carry over for $\handwritten{\underline}{C}$ exactly as for $\handwritten{\underline}{C}^{\handwritten{0}}$.

\handwritten{\uline}{References}.

1. A.O.L.Atkin,Proof of a conjecture of Ramanujan,Glasgow Math.J. \handwritten{\uwave}{8}(1967),14-32.

2. A.O.L.Atkin,Congruences for modular forms,Proceedings of IBM Conference on computers in mathematical research,Blaricum 1966,published North Holland(1967).

3.A.O.L.Atkin and J.N.O'Brien,Some properties of $p(n)$ and $c(n)$ modulo powers of 13,Trans. Amer. Math. Soc. \handwritten{\uwave}{126}(1967),442-459.

4. R.C.Gunning,Lectures on modular forms,Princeton 1962.

5. J.Lehner,Divisibility properties of the Fourier coefficients of the modular invariant $j(\handwritten{\underline}{\handwritten{\tau}})$,Amer. J.Math.\handwritten{\uwave}{71}(1949),\deleted{\unsure{1}\illegible{1}\unsure{5|3}}\inserted{\handwritten{136-148}}.\reversemarginpar\margin{\handwritten{\unsure{Fin|Fiu}}}\normalmarginpar

6. J.Lehner,Further congruence properties of the Fourier coefficients of the modular invariant $j(\handwritten{\underline}{\handwritten{\tau}})$,Amer.J.Math.\handwritten{\uwave}{71}(1949),373-386.

7.J.Lehner,A short course in automorphic functions,New York 1965.

8. R.A.Rankin,Elliptic \deleted{\unsure{M|m}}\handwritten{m}odular \deleted{\unsure{F|f}}\handwritten{f}unctions and \deleted{\unsure{F|f}}\handwritten{f}orms,Lecture notes ,Bloomington,Indiana,1964.

\deleted{\unsure{8|0}}9. R.A.Rankin and J.M.Rushforth,The coefficients of certain integral modular forms,Proc. Cambridge Phil.Soc. \handwritten{\uwave}{50}(1954),\unsure{305|395}-\unsure{308|300}.\margin{\handwritten{\#}}

\unsure{The} \unsure{Atlas} \unsure{Computer} Laboratory,\unsure{Chilton|Chillton},Didcot,England\unsure{:|.}
% archive sheet number: 63

